\section{Individual spin classes and their intersections}\label{sec:individual}
On covers with two simple branch points the residue description determines a class for each spin structure. Its character decomposition gives a lattice indexed by the two-torsion of the Jacobian. On covers with four simple branch points there is a further distinction: the Euler integral of two spins has contributions on several components of the stable-map stack. We isolate the component of finite smooth covers and compute the ten individual intersection pairings on its closure in the genus-two cubic case.
\subsection{The lattice of spin Euler classes}\label{p8-sec:lattice}
The individual Euler classes on the Hurwitz surface contain
information beyond their pairwise integrals. Fix a spin line
$\theta_0$, put $G=H^1(C;\ZZ/2)$, and use the flat sign line to
write $\theta_0 z$ for its twist by $z\in G$. On the complement of
a node's image use the permutation action of the normalization.
At a triple point its local monodromy is even. In either case the
permutation sign on the complement of the collision extends to a
class $\alpha\in G$. For a boundary divisor $B$ in
\eqref{p8-eq:intro-divisor} write this class as $\alpha_B$. It is
constant on $B$: locally the covering data away from a disk about
the moving point do not vary, and removing and filling the point
induces an isomorphism on the extended sign characters.
Define
\[
 D_\alpha=\sum_{B:\alpha_B=\alpha} s_{\theta_0}(B)[B],
 \quad
 \Lambda=\sum_{\alpha\in G}\ZZ D_\alpha,\qquad
 \Lambda_{\rm spin}=\sum_{z\in G} \ZZ\cdot4e(V_{\theta_0 z}),
\]
as lattices in $H^2(\mathcal F_m;\QQ)$. Here $V_\theta$ is the
oriented spin cokernel.
\begin{theorem}\label{p8-thm:lattice}
For every $m\ge2$, the $D_\alpha$ are independent, the spin Euler
classes span the same rational space, and
\[
     \Lambda/\Lambda_{\rm spin}
          \simeq \bigoplus_{j=1}^{2g}
               (\ZZ/2^j\ZZ)^{\binom{2g}{j}}.
\]
The least positive integer taking any one of the $e(V_\theta)$
into $\Lambda$ is four.
\end{theorem}
\begin{proof}
Every $\alpha$ occurs already on a nodal divisor. For nonzero
$\alpha$, choose surface generators with value one on $a_1$ alone.
If $m$ is even, send $a_1$ to an $m$-cycle. If $m$ is odd, send
$a_1$ to a transposition and $a_2$ to an $m$-cycle. The other chosen
generators are trivial, so this gives a connected unramified cover
with sign $\alpha$. For $\alpha=0$ and odd $m$ an $m$-cycle suffices.
For even $m\ge4$, assign an $(m-1)$-cycle and a $3$-cycle involving
one of its letters and the last sheet so that the generated action
is transitive; both permutations are even. In each instance glue a
pair of points in a fiber. For $m=2$ and $\alpha=0$ take two trivial
sheets and glue their points over the chosen value. Distinct reduced
boundary divisors are disjoint with strictly negative self-intersection.
The $D_\alpha$ therefore have mutually orthogonal nonzero squares.

Set $q_{\theta_0}(z)=h^0(\theta_0 z)-h^0(\theta_0)\bmod2$.
The ratio of the two boundary parities, by the orthogonal trace-bundle
calculation in Section~\ref{p8-sec:characters}, gives
\[
 -4e(V_{\theta_0 z})
    =(-1)^{m q_{\theta_0}(z)}
       \sum_{\alpha\in G}
          (-1)^{\langle z\smile\alpha,[C]\rangle}D_\alpha .
\]
The common row signs do not change the generated lattice. In dual
binary bases for the nondegenerate cup product the character matrix
is the $2g$-fold tensor power of the matrix with rows $(1,1)$ and
$(1,-1)$. Integral row and column operations diagonalize this
two-by-two matrix with entries $1,2$. Tensoring those unimodular
operations gives the Smith factors $2^j$ with the stated
multiplicities. Every coefficient of a single $e(V_\theta)$ in the
independent $D_\alpha$ is $\pm1/4$, proving the last assertion.
\end{proof}
The theorem compares two lattices generated by geometric classes in
rational cohomology. It does not require that the individual
boundary divisors span the integral cohomology of the stack.
\subsection{Four simple branch points}\label{sec:excess}
Fix two theta characteristics \(a,b\) on \(C\). Write \(c_d\) for the signed connected étale degree and \(F_d\) for the finite-map Euler integral on the first Hurwitz surface. Let \(M_d\) be the closure of the smooth connected covers with four distinct simple branch points, with the sum of its reduced four-dimensional component cycles. Write \(V_a,V_b\) for the real spin cokernels and put \(I_d=\int_{[M_d]}e(V_a)e(V_b)\).
\begin{proposition}\label{prop:excess}
For every \(d\ge2\), putting \(t=d(g-1)\), one has
\[
 L_{d,H_d+2}=I_d-\frac{t}{12}F_d+\frac{t(5t-1)}{1440}c_d.
\]
For a genus-two target and distinct even spins this gives \(I_2=0,\ L_{2,5}=-1/160,\ n_{2,5}=0\) and \(I_3=3/2\). For a repeated even spin on the same target, \(I_3=161/6\).
\end{proposition}
\begin{proof}
First remove the open and closed summands consisting of maps of degree one over a connected étale target cover of full degree \(d\); its lift is taken modulo its deck group. Their degree-one series contributes \([v^{2t+4}]x^t c_d=t(5t-1)c_d/1440\). Covers factoring through a smaller étale target remain. Normalize a map in the complement. The sum of half the ramification degree, the genera of contracted normalizations and the excess from the dual graph is two. The remaining positive-dimensional contracted strata have a pointed elliptic tail on a two-branch finite core, an elliptic component bridging two sheets, or a rational tripod. The first family has a proper parametrization by the universal pointed first Hurwitz surface times \(\overline{\mathcal M}_{1,1}\), including the inserted bridge when the mark meets a node. The degree on the generic image is one.

The other strata lie on \(M_d\): attaching a degree-two bubble by the pencil \(|p+q|\) smooths an elliptic bridge with attaching pair \(p,q\). For the tripod attach the rational function \(z+z^{-1}+(z-1)^{-1}\) at its three simple poles. It has four simple branch values, since the binary discriminant in the target coordinate is \(w^4-2w^3-5w^2+6w-23\), of discriminant \(-5619712\). At a direct node glue rational bubble covers of the matching contact orders and smooth the balanced equations after base change. In each case the stabilized target remains \(C\). The remaining support is the union of \(M_d\) and the elliptic-tail image, both pure of dimension four. A smooth simply branched cover has generic virtual multiplicity one, and at a generic tail the derivative at the unramified attaching point identifies the smoothing line with \(T_{f(p)}C\otimes H^1(E,\mathcal O_E)\). The generic virtual multiplicity there is also one. A cycle of the top dimension is determined by these multiplicities. On the tail the matching constants cancel, leaving \(\mathbb E^\vee\otimes\operatorname{ev}^*(a\oplus b)\) with top class \(-\lambda\operatorname{ev}^*c_1(K_C)\). Its integral along the two new parameters is \(-t/12\). This proves the formula; Calculation~B3 of \cite{WueSupplement} gives the full normalization along the lower-dimensional strata. Substitution of the degree-two and cubic coefficients of Theorem~\ref{thm:spin} and the signed finite sums of Sections~\ref{p8-sec:hurwitz}--\ref{p8-sec:characters} gives the stated values; for the repeated spin one uses the orientable balanced series. See also Calculations~B1--B4 of \cite{WueSupplement}.
\end{proof}

For genus two, distinct even spins give the entire connected double series
\[
 L_{2,k}=(-1)^k\frac{2^{2k-2}-4}{(2k-2)!},\qquad k\ge3 .
\]
Indeed the disconnected coefficient of \(Q^2\) vanishes and its connected coefficient is \(-x^2/2\). Subtracting the balanced factor in degree one leaves no elementary double exponent at genus five or higher. A separate geometric bound for reduced multisections is available. If the normal equations are independent line equations of degrees at most \(d_{\max}\), and the reduced graph is generically one-to-one without contracted components, the ramification degree plus twice the number of identified nodes is at most \(m(m-1)d_{\max}\). On the normalized ordered-pair cover the difference of the two solutions satisfies a real line equation with positive zero orders. Its zeros at ramification and at the ordered pairs identified by a node give the inequality by integrating the first Chern class. Pairs on which one component difference is identically zero are counted using the other line. In the triangular case one first separates distinct scalar projections and then applies the line bound in each equal-projection group. The projected-jet variation, the prime-degree consequence and the restriction at composite towers are detailed in Calculations~A1--A10 of \cite{WueSupplement}. In particular for noneffective spins and prime degree \(p\) the elementary genus is at most \(1+(g-1)\binom{p+1}{2}\).

\subsection{The ten even spins in genus two}
On \(M_3\subset\overline{\mathcal M}_6(C,3)\) put \(W_a=-R\pi_*f^*a\) and \(\Delta_a=(c_2(W_a)-e(V_a))/2\).
\begin{theorem}\label{thm:gram}
For every smooth genus-two target the ten even spins have \(\int_{[M_3]}\Delta_a^2=6\) and \(\int_{[M_3]}\Delta_a\Delta_b=-1/3\) for \(a\ne b\). The ten classes span a positive-definite space of rank ten.
\end{theorem}
\begin{proof}
Let \(\iota\) be the hyperelliptic involution of \(C\). For distinct even spins \(a,b\), the two residue-one Szeg\H o kernels have the same polar part along the diagonal. Their squared difference restricts there to a quadratic differential \(Q_{ab}\). The common section locus on the genus-five finite Hurwitz surface is its zero scheme. To see the condition, consider the trace-free part of \(f_*\mathcal O_Y\). The two pulled-back sections identify it with an elementary enlargement of \(a^{-1}\oplus b^{-1}\) at a point \(p\). In square frames at \(p\), the normalized meromorphic sections of the two spins have expansions \(z^{-1}+c_a z+\cdots\) and \(z^{-1}+c_b z+\cdots\). Expressing the cubic multiplication in the enlarged frame leaves one regularity condition, \(c_a=c_b\), or \(Q_{ab}(p)=0\). This calculation works over local Artin bases. The resulting cover is unique, smooth, with one triple ramification point over \(\iota(p)\); its monodromy is \(S_3\), so it has no deck factor. The full cubic multiplication and the boundary classification are in Calculation~C1 and Calculation~C2, \S\S2--3 of \cite{WueSupplement}. On the four-cycle \(M_3\) the common support consists of the disjoint curves \(E_p\simeq\overline{\mathcal M}_{1,1}\) obtained by attaching an elliptic tail at that triple point.

Suppose that \(p\) is a simple zero of \(Q_{ab}\). At the triple point the map deforms as \(w^3+tw+v\); let \(w=x\) mark the attachment and let \(\tau\) smooth the tail node. The versal equation is \(\tau B=0\), where \(B|_{\tau=0}=t+3x^2\) up to a unit, by the principal-part obstruction to smoothing. Thus \(B=0\) is a smooth chart of the main component and \(\tau=0\) cuts its tail divisor. On this divisor the ramification marking satisfies \(t=-3x^2\). Each supported correction \(\Delta_\theta\) is the class of its kernel surface with coefficient one, by complex-to-real determinant continuation in Calculation~C2, \S5 of \cite{WueSupplement}. On the finite surface the two kernel divisors meet with length one at this simple zero. Pullback under the degree-two ramification marking makes their intersection inside the tail divisor equal to \(2[E_p]\). Its normal line in the main component is the node-smoothing line, of first Chern class \(-\lambda\) on \(E_p\). The self-intersection formula for that divisor therefore gives
\[
 \operatorname{Contr}_{E_p}(\Delta_a\Delta_b)
     =-2\int_{\overline{\mathcal M}_{1,1}}\lambda=-\frac1{12}.
\]
The stack Hodge degree includes the pointed-tail automorphisms. Since \(K_C^2\) has degree four, four simple zeros give total mixed degree \(-1/3\).

Put \(S_a=\int\Delta_a^2\). The equal-spin integral is \(161/6\) from Proposition~\ref{prop:excess}; the distinct-spin integral is \(3/2\). The Chern class of the complex index is independent of \(a\) rationally, since spin differences are torsion. Subtracting one half the square of the difference of the two real Euler classes from one half the sum of their squares gives
\[
    \frac32=\frac{161}{6}-2(S_a+S_b)+4\int\Delta_a\Delta_b .
\]
It follows that \(S_a+S_b=12\) whenever \(Q_{ab}\) has four simple zeros.

Write the six hyperelliptic branch values as \(e_i\). Each spin is an unordered partition into two triples. Given two distinct partitions choose sides sharing two elements with monic root polynomial \(P\); the opposite sides share two with root polynomial \(R\). Call the exchanged values \(c,d\). The hyperelliptic Szeg\H o formula \cite[Appendix~C, (C.3)--(C.4)]{dhs} gives \(Q_{ab}=(c-d)(P'R-PR')(dx/y)^2/4\). The binary quadratic \(P'R-PR'\) has discriminant \(4\prod_{P(a_i)=R(b_j)=0}(a_i-b_j)\ne0\). Only \(c,d\) can be exceptional branch roots. At a fixed \(c\) this happens precisely when \(u_{a_1}+u_{a_2}=u_{b_1}+u_{b_2}\) for \(u_e=1/(c-e)\). Of five distinct \(u_e\) there are at most three such disjoint-pair equalities. Indeed project to distinct reals and let their successive gaps be \(A,B,C',D>0\). The five possible equations are \(A=C'\), \(A=C'+D\), \(A=D\), \(A+B=D\), \(B=D\); at most three coexist. The graph on ten spins with the four-simple-zero pairs as edges has at least \(45-18=27\) edges. At a fixed spin and branch value at most one omitted element realizes the equality, so every vertex has degree at least three. A disconnected graph with this minimum degree would have at most \(21\) edges, a bipartite graph at most \(25\). Connectedness and an odd cycle force \(S_a=6\) throughout. The displayed equation now gives \(-1/3\) for every distinct pair. The Gram matrix has eigenvalues \(3\) and \(19/3\), the latter with multiplicity nine.
\end{proof}

The local contribution at a zero of order \(j\) is \(-j/12\), including a nonsimple zero. Indeed the supported Euler products on the different closed elliptic-tail curves add by excision. At each simple zero the local calculation in the proof applies irrespective of the other zeros. One double zero must therefore contribute \(-1/6\) by the global mixed pairing and the two simple contributions. With two double zeros send their underlying branch values to \(0,\infty\). They are the ramification points of the map \(P/R\); its deck involution is \(z\mapsto-z\). Thus the common pairs are \(\{\pm u\}\), \(\{\pm v\}\). Writing \(k=uv\), the hyperelliptic curve has equation \(y^2=z(z^2-u^2)(z^2-v^2)\) and automorphism \((z,y)\mapsto(k/z,k^{3/2}y/z^3)\). It exchanges the two supports and preserves each spin because it sends its triple to its complement. Choose square-compatible lifts and invariant metrics by averaging over the finite extension by signs. The spin orientation continued from the complex determinant and the supported Euler product are then preserved. Hence the two contributions agree, each being \(-1/6\).

On \(y^2=\prod_{j=1}^6(x-j)\), the two exceptional pairs, indexing a spin by its triple containing \(1\), are \(126,135\) and \(135,156\). The smooth-source kernel pairing equals \(5\) for every even spin: the ramification-tail term equals \(1\) and the complete square is \(6\). Calculations~C3--C4 of \cite{WueSupplement} give the characteristic-degree calculation on the singular spectral incidence and its separate tail calculation. Identifying the two compactifications at cusps and tacnodes remains open.
